%!TeX root=../persuasiontop.tex
\chapter[Chapter \thechapter]{}
	
	\begin{figure}[t!]
\centering
\includegraphics[width=\linewidth]{nouveau04}
\end{figure}

\lettrine[lraise=0.3]{H}{e} was not Mr Wentworth, the former curate of Monkford, however suspicious appearances may be, but a Captain Frederick Wentworth, his brother, who being made commander in consequence of the action off St Domingo, and not immediately employed, had come into Somersetshire, in the summer of 1806; and having no parent living, found a home for half a year at Monkford. He was, at that time, a remarkably fine young man, with a great deal of intelligence, spirit, and brilliancy; and Anne an extremely pretty girl, with gentleness, modesty, taste, and feeling. Half the sum of attraction, on either side, might have been enough, for he had nothing to do, and she had hardly anybody to love; but the encounter of such lavish recommendations could not fail. They were gradually acquainted, and when acquainted, rapidly and deeply in love. It would be difficult to say which had seen highest perfection in the other, or which had been the happiest: she, in receiving his declarations and proposals, or he in having them accepted.

A short period of exquisite felicity followed, and but a short one. Troubles soon arose. Sir Walter, on being applied to, without actually withholding his consent, or saying it should never be, gave it all the negative of great astonishment, great coldness, great silence, and a professed resolution of doing nothing for his daughter. He thought it a very degrading alliance; and Lady Russell, though with more tempered and pardonable pride, received it as a most unfortunate one.

Anne Elliot, with all her claims of birth, beauty, and mind, to throw herself away at nineteen; involve herself at nineteen in an engagement with a young man, who had nothing but himself to recommend him, and no hopes of attaining affluence, but in the chances of a most uncertain profession, and no connexions to secure even his farther rise in the profession, would be, indeed, a throwing away, which she grieved to think of! Anne Elliot, so young; known to so few, to be snatched off by a stranger without alliance or fortune; or rather sunk by him into a state of most wearing, anxious, youth-killing dependence! It must not be, if by any fair interference of friendship, any representations from one who had almost a mother's love, and mother's rights, it would be prevented.

Captain Wentworth had no fortune. He had been lucky in his profession; but spending freely, what had come freely, had realized nothing. But he was confident that he should soon be rich: full of life and ardour, he knew that he should soon have a ship, and soon be on a station that would lead to everything he wanted. He had always been lucky; he knew he should be so still. Such confidence, powerful in its own warmth, and bewitching in the wit which often expressed it, must have been enough for Anne; but Lady Russell saw it very differently. His sanguine temper, and fearlessness of mind, operated very differently on her. She saw in it but an aggravation of the evil. It only added a dangerous character to himself. He was brilliant, he was headstrong. Lady Russell had little taste for wit, and of anything approaching to imprudence a horror. She deprecated the connexion in every light.

Such opposition, as these feelings produced, was more than Anne could combat. Young and gentle as she was, it might yet have been possible to withstand her father's ill-will, though unsoftened by one kind word or look on the part of her sister; but Lady Russell, whom she had always loved and relied on, could not, with such steadiness of opinion, and such tenderness of manner, be continually advising her in vain. She was persuaded to believe the engagement a wrong thing: indiscreet, improper, hardly capable of success, and not deserving it. But it was not a merely selfish caution, under which she acted, in putting an end to it. Had she not imagined herself consulting his good, even more than her own, she could hardly have given him up. The belief of being prudent, and self-denying, principally for \textit{his} advantage, was her chief consolation, under the misery of a parting, a final parting; and every consolation was required, for she had to encounter all the additional pain of opinions, on his side, totally unconvinced and unbending, and of his feeling himself ill used by so forced a relinquishment. He had left the country in consequence.

A few months had seen the beginning and the end of their acquaintance; but not with a few months ended Anne's share of suffering from it. Her attachment and regrets had, for a long time, clouded every enjoyment of youth, and an early loss of bloom and spirits had been their lasting effect.

More than seven years were gone since this little history of sorrowful interest had reached its close; and time had softened down much, perhaps nearly all of peculiar attachment to him, but she had been too dependent on time alone; no aid had been given in change of place (except in one visit to Bath soon after the rupture), or in any novelty or enlargement of society. No one had ever come within the Kellynch circle, who could bear a comparison with Frederick Wentworth, as he stood in her memory. No second attachment, the only thoroughly natural, happy, and sufficient cure, at her time of life, had been possible to the nice tone of her mind, the fastidiousness of her taste, in the small limits of the society around them. She had been solicited, when about two-and-twenty, to change her name, by the young man, who not long afterwards found a more willing mind in her younger sister; and Lady Russell had lamented her refusal; for Charles Musgrove was the eldest son of a man, whose landed property and general importance were second in that country, only to Sir Walter's, and of good character and appearance; and however Lady Russell might have asked yet for something more, while Anne was nineteen, she would have rejoiced to see her at twenty-two so respectably removed from the partialities and injustice of her father's house, and settled so permanently near herself. But in this case, Anne had left nothing for advice to do; and though Lady Russell, as satisfied as ever with her own discretion, never wished the past undone, she began now to have the anxiety which borders on hopelessness for Anne's being tempted, by some man of talents and independence, to enter a state for which she held her to be peculiarly fitted by her warm affections and domestic habits.

They knew not each other's opinion, either its constancy or its change, on the one leading point of Anne's conduct, for the subject was never alluded to; but Anne, at seven-and-twenty, thought very differently from what she had been made to think at nineteen. She did not blame Lady Russell, she did not blame herself for having been guided by her; but she felt that were any young person, in similar circumstances, to apply to her for counsel, they would never receive any of such certain immediate wretchedness, such uncertain future good. She was persuaded that under every disadvantage of disapprobation at home, and every anxiety attending his profession, all their probable fears, delays, and disappointments, she should yet have been a happier woman in maintaining the engagement, than she had been in the sacrifice of it; and this, she fully believed, had the usual share, had even more than the usual share of all such solicitudes and suspense been theirs, without reference to the actual results of their case, which, as it happened, would have bestowed earlier prosperity than could be reasonably calculated on. All his sanguine expectations, all his confidence had been justified. His genius and ardour had seemed to foresee and to command his prosperous path. He had, very soon after their engagement ceased, got employ: and all that he had told her would follow, had taken place. He had distinguished himself, and early gained the other step in rank, and must now, by successive captures, have made a handsome fortune. She had only navy lists and newspapers for her authority, but she could not doubt his being rich; and, in favour of his constancy, she had no reason to believe him married.

How eloquent could Anne Elliot have been! how eloquent, at least, were her wishes on the side of early warm attachment, and a cheerful confidence in futurity, against that over-anxious caution which seems to insult exertion and distrust Providence! She had been forced into prudence in her youth, she learned romance as she grew older: the natural sequel of an unnatural beginning.

With all these circumstances, recollections and feelings, she could not hear that Captain Wentworth's sister was likely to live at Kellynch without a revival of former pain; and many a stroll, and many a sigh, were necessary to dispel the agitation of the idea. She often told herself it was folly, before she could harden her nerves sufficiently to feel the continual discussion of the Crofts and their business no evil. She was assisted, however, by that perfect indifference and apparent unconsciousness, among the only three of her own friends in the secret of the past, which seemed almost to deny any recollection of it. She could do justice to the superiority of Lady Russell's motives in this, over those of her father and Elizabeth; she could honour all the better feelings of her calmness; but the general air of oblivion among them was highly important from whatever it sprung; and in the event of Admiral Croft's really taking Kellynch Hall, she rejoiced anew over the conviction which had always been most grateful to her, of the past being known to those three only among her connexions, by whom no syllable, she believed, would ever be whispered, and in the trust that among his, the brother only with whom he had been residing, had received any information of their short-lived engagement. That brother had been long removed from the country and being a sensible man, and, moreover, a single man at the time, she had a fond dependence on no human creature's having heard of it from him.

The sister, Mrs Croft, had then been out of England, accompanying her husband on a foreign station, and her own sister, Mary, had been at school while it all occurred; and never admitted by the pride of some, and the delicacy of others, to the smallest knowledge of it afterwards.

With these supports, she hoped that the acquaintance between herself and the Crofts, which, with Lady Russell, still resident in Kellynch, and Mary fixed only three miles off, must be anticipated, need not involve any particular awkwardness.